% Amélioration de l'affichage des figures et tableaux pour une présentation professionnelle
\usepackage{float}             % Pour un meilleur contrôle des positions des figures
\usepackage[export]{adjustbox} % Pour ajouter des cadres et ajuster les images
\usepackage{subcaption}        % Pour les sous-figures
\usepackage{wrapfig}           % Pour enrouler le texte autour des figures
\usepackage{booktabs}          % Pour des tableaux plus professionnels
\usepackage{multirow}          % Pour fusionner des lignes dans les tableaux
\usepackage{colortbl}          % Pour colorer les tableaux

% Les couleurs sont déjà définies dans color-styling.tex, nous les utilisons ici

% Configuration avancée des légendes
\usepackage{caption}
\captionsetup{
  font=small,                % Taille de police pour les légendes
  labelfont={bf,color=captionColor}, % Label en gras et coloré
  textfont=it,               % Texte en italique
  width=0.9\textwidth,       % Largeur de la légende
  format=hang,               % Format suspendu
  skip=10pt,                 % Espace entre figure et légende
  margin=10pt,               % Marge
  justification=justified    % Justification du texte
}

% Définir les styles de cadre pour les figures
\newcommand{\figurestyle}[1]{%
  \fbox{\includegraphics[width=#1\textwidth, keepaspectratio]}%
}

% Commande pour les figures avec cadre simple
\newcommand{\framedimage}[3][0.8]{%
  \begin{figure}[H]
    \centering
    \fcolorbox{figBorder}{white}{\includegraphics[width=#1\textwidth, keepaspectratio]{#2}}
    \caption{#3}
  \end{figure}
}

% Commande pour les figures avec ombre
\newcommand{\shadowimage}[3][0.8]{%
  \begin{figure}[H]
    \centering
    \includegraphics[width=#1\textwidth, keepaspectratio, frame, shadow]{#2}
    \caption{#3}
  \end{figure}
}

% Commande pour des figures côte à côte
\newcommand{\sidebysideimages}[6][0.45]{%
  \begin{figure}[H]
    \centering
    \begin{subfigure}[b]{#1\textwidth}
      \centering
      \includegraphics[width=\textwidth, keepaspectratio]{#2}
      \caption{#3}
    \end{subfigure}%
    \hfill
    \begin{subfigure}[b]{#1\textwidth}
      \centering
      \includegraphics[width=\textwidth, keepaspectratio]{#4}
      \caption{#5}
    \end{subfigure}
    \caption{#6}
  \end{figure}
}
